%% =====================================================================
%% INTRODUCTION -- MACRO ROADMAP (whole-section level)
%% ---------------------------------------------------------------------
%% What this Introduction must convey to the reader:
%%   A body-referenced multisensory representation -- peripersonal
%%   space (PPS) -- is not fixed. It is continually re-drawn by the
%%   brain's ongoing estimate of body state. Interoceptive signals
%%   belong in that estimate, but the route by which they reach the
%%   near-body sigmoid is not specified. Respiration is the best
%%   candidate route: it is the most controllable and most broadly
%%   influential body rhythm, it organises perception and multisensory
%%   integration across modalities, and it can be coupled to an
%%   external visual signal to make the internal breath state
%%   exteroceptively salient. The predicted intersection -- respiratory
%%   phase x PPS, augmented by breath-coupled visual feedback -- has
%%   never been directly tested. The present study runs that test.
%%
%% ARGUMENTATIVE CHAIN (subsection-level, one arrow per step):
%%   [Opener]     PPS is a plastic, body-referenced MSI interface
%%                indexed by a sigmoid (extent, slope).
%%        -->
%%   [Sub 1]      PPS is anchored to the brain's estimate of body
%%                state, not the physical body.
%%        -->
%%   [Sub 2]      Interoceptive signals modulate the sigmoid; PPS
%%                is an intero-exteroceptive interface -- but the
%%                channel is unspecified.
%%        -->
%%   [Sub 3]      Respiration is the candidate channel: it organises
%%                perception and, critically, multisensory integration
%%                itself.
%%        -->
%%   [Sub 4]      Yet the direct respiratory-phase x PPS test is
%%                absent, AND the analogous plasticity manipulation
%%                (breath-coupled visual environment) has not been
%%                tried.
%%        -->
%%   [Sub 5]      The present design tests both simultaneously:
%%                phase-locked PPS under Viscereality vs Control;
%%                four pre-registered hypotheses.
%% =====================================================================

\section{Introduction}

%% ---------------------------------------------------------------------
%% OPENER PARAGRAPH
%% Goal (one idea): Introduce PPS as the body's multisensory interface
%% with the world, quantified by a sigmoid with two parameters
%% (extent, sharpness) that together index where and how crisply the
%% brain draws the self-world boundary. This frames the whole paper's
%% outcome measure.
%% Role in chain: establishes the object of study.
%% ---------------------------------------------------------------------
Embodied cognition takes the mind to be built on, and continually shaped by, the sensorimotor exchanges between a body and its environment \citep{Glenberg2010, Barsalou2008, barsalou2008grounded}. A privileged zone of that exchange is \emph{peripersonal space} (PPS), the multisensory region immediately surrounding the body within which incoming signals are integrated and prepared for action \citep{rizzolatti1997space, diPellegrino2015, Serino2019}. PPS is typically operationalised through audio--tactile or visuo--tactile paradigms in which a tactile stimulus is delivered while an auditory or visual stimulus approaches or recedes from the body; reaction times form a sigmoid over proximity, yielding two parameters of theoretical interest: the \emph{extent} of the near-body zone (its central point) and its \emph{sharpness} (the slope at the central point) \citep{Canzoneri2012, teneggi2013social, Bufacchi2018, Basile2024}. On current accounts, PPS is not a distinct cortical module but a case of multisensory integration operating under body-centred coordinate frames and an additional proximity-gating term \citep{noel_2018, bertoni_2025, straka_2022}. It therefore offers a behavioural handle on how the brain continually re-draws the boundary between self and world.

\subsection{Peripersonal space is plastic and tracks the state of the body}

%% ---------------------------------------------------------------------
%% SUBSECTION 1 GOAL
%% What this subsection conveys: every PPS plasticity finding to date
%% can be read as a shift in the brain's multisensorially estimated
%% body state, not a shift in the physical body. This is the
%% theoretical foothold that lets the argument move in the next
%% subsection from exteroceptive to interoceptive modulators: if
%% PPS is anchored to a *body estimate*, interoceptive signals (which
%% feed that estimate) should also shift PPS.
%% ---------------------------------------------------------------------

%% ---
%% PARAGRAPH GOAL (one idea): Tool use, social, body-illusion, and
%% vestibular manipulations all shift PPS in a way consistent with
%% a single underlying theme -- PPS tracks the brain's body estimate,
%% not the physical body. This sets up the "what else feeds the
%% body estimate?" question.
%% ---
The sigmoid expressed by PPS is remarkably malleable. Tool use extends it to incorporate the tool \citep{iriki1996coding, maravita2004tools}; cooperative social interaction merges self and partner space \citep{teneggi2013social}; body-swap and full-body illusions relocate PPS with the illusorily perceived self-location \citep{noel2015peripersonal, noel2015fullbody}; and vestibular stimulation reshapes it in line with altered spatial orientation \citep{pfeiffer2018vestibular}. Across these manipulations a single theme recurs: PPS is anchored to the multisensorially estimated state of the body, not to the physical body per se \citep{blanke2015behavioral, Grivaz2017, serino2015bodypart}. Whatever shifts the brain's ongoing estimate of ``where the body is'' accordingly shifts the near-body coordinate frame within which incoming signals are integrated.

\subsection{Interoception and the softening of the self--world boundary}

%% ---------------------------------------------------------------------
%% SUBSECTION 2 GOAL
%% What this subsection conveys: interoceptive signals also modulate
%% PPS, establishing PPS as an intero-exteroceptive interface. But the
%% route by which internal signals reach the sigmoid is unspecified --
%% setting up respiration as the candidate in Subsection 3.
%% ---------------------------------------------------------------------

%% ---
%% PARAGRAPH GOAL (one idea): Cardiac interoceptive accuracy narrows
%% PPS; brief meditation softens it; mindfulness and sensory-attenuation
%% interventions dissolve body boundaries. Together these place PPS at
%% an intero-exteroceptive interface, but the mechanistic question
%% (which interoceptive signal, through which temporal channel?) is
%% unresolved.
%% ---
A more recent and smaller literature extends this logic from exteroceptive to interoceptive signals. In a foundational study, \citet{Ardizzi2018} showed that individuals with higher cardiac interoceptive accuracy have \emph{narrower} PPS boundaries, with the effect moderated by a dispositional tendency to attend to internal sensations. Brief focused-attention meditation softens the sharpness of the PPS sigmoid without contracting its extent \citep{Chiarella2024}; guided meditation abolishes the near--far attentional enhancement of auditory evoked responses in non-meditators \citep{nguyen_2020}; mindfulness training relaxes perceived body boundaries and favours allocentric spatial frames \citep{hanley_2020}; and sensory-attenuation contexts such as Floatation-REST produce body-boundary dissolution that mediates anxiety reduction \citep{hruby_2024}. Read together with the broader case for breath, heart, and gut as shapers of self-experience \citep{seth2013interoceptive, barrett2016theory, Park2019, bogdanova_2021}, these findings converge on a view of PPS as an \emph{intero--exteroceptive interface}: the zone in which interoceptive priors about the body meet the statistics of the surrounding world. What remains underspecified is \emph{which} interoceptive signals reach the PPS computation, and through which temporal channel.

\subsection{Respiration as an organiser of perception and multisensory integration}

%% ---------------------------------------------------------------------
%% SUBSECTION 3 GOAL
%% What this subsection conveys: respiration is the most plausible
%% interoceptive channel for reaching PPS. It is uniquely available
%% (both voluntary and automatic, rich afferents), it modulates
%% perception across modalities, and -- critically -- it modulates
%% multisensory integration itself, which is the direct bridge to PPS.
%%
%% Two paragraphs, two ideas:
%%   (a) respiration is a broad perceptual organiser
%%   (b) respiration modulates MSI specifically (Saltafossi), which
%%       places it in a predictive-coding framework as a precision-
%%       allocation scaffold.
%% ---------------------------------------------------------------------

%% ---
%% PARAGRAPH 3a GOAL (one idea): respiration, among body rhythms, is
%% uniquely positioned and empirically shown to organise cortical
%% excitability and perception across modalities. Establish the
%% biological plausibility of respiration-as-interoceptive-channel
%% for PPS.
%% ---
Breath is the most accessible candidate. Unlike other autonomic rhythms, respiration is both voluntarily controllable and automatically driven, and it carries chemoreceptor, mechanoreceptor, and proprioceptive afferents that converge on brainstem, insular, and limbic targets \citep{Chan2024, Varga2017, heck2017breathing, Zelano2022}. A decade of work has established that respiratory phase modulates cortical excitability and perception well beyond the olfactory system. Nasal inhalation entrains oscillations in piriform cortex, amygdala, and hippocampus, and correspondingly speeds fearful-face discrimination and episodic retrieval \citep{zelano2016nasal, mizuhara_2023}. Respiration is coupled to posterior alpha power and to visual detection sensitivity, with excitability changes leading performance fluctuations \citep{kluger_2021}. Tactile detection peaks in the first quadrant after expiration onset, and participants actively phase-lock their breathing to expected stimuli \citep{grund_2021}. Breath--heart coupling via respiratory sinus arrhythmia shapes heartbeat-evoked potentials and the gating of somatosensory input \citep{Zaccaro2022, Heck2022}.

%% ---
%% PARAGRAPH 3b GOAL (one idea): Saltafossi 2025 is the direct bridge
%% from respiratory-phase effects to MSI, and therefore the proximate
%% empirical anchor for predicting respiratory-phase effects on PPS
%% (since PPS is MSI under body-centred frames). Place this within a
%% predictive-coding account of respiration as a precision scaffold
%% so the prediction is principled, not just empirical.
%% ---
Most consequentially for the present work, \citet{saltafossi_2025} showed that \emph{respiratory phase modulates multisensory integration itself}: using the race-model inequality, they found that audio--tactile and audio--visual integration peaks at the expiration-to-inspiration transition, dissociating multisensory gain from unimodal speed. Predictive-coding accounts formalise this by casting respiration as a temporal scaffold along which the brain allocates precision to competing sources of evidence about the body and the surrounding world, with breath-phase modulating local and global neural gain across interoceptive and exteroceptive channels \citep{allen_2021, braendholt_2023, engelen_2023}. On this view, breath is not an arousal confound but an organising rhythm for sensory inference. Because PPS is multisensory integration under body-centred reference frames and a proximity-gating term \citep{noel_2018, bertoni_2025}, the same respiratory-phase organisation should, in principle, propagate to the PPS sigmoid.

\subsection{A missing test and a missing manipulation}

%% ---------------------------------------------------------------------
%% SUBSECTION 4 GOAL
%% What this subsection conveys: two convergent literatures have
%% arrived at the same question without testing it, and the
%% corresponding plasticity manipulation has not been attempted.
%% This is the explicit gap that the present study fills.
%%
%% Two paragraphs, two ideas:
%%   (a) empirical gap: no respiratory-phase-locked PPS measurement
%%   (b) paradigmatic gap: no breath-coupled environmental-synchrony
%%       manipulation of PPS; cardio-visual precedent shows the
%%       mechanism is available.
%% ---------------------------------------------------------------------

%% ---
%% PARAGRAPH 4a GOAL (one idea): the direct test of respiratory phase
%% on the PPS sigmoid has not been run; the one near-precedent
%% (Spaccasassi) used breathing as anxiety induction, not as a phase
%% manipulation. State the empirical gap plainly so the paper's
%% contribution is clear.
%% ---
Two literatures have therefore converged on the same question without directly addressing it. PPS is a body-referenced multisensory integration process shaped by interoceptive signals \citep{Ardizzi2018, bogdanova_2021}; respiratory phase is a powerful organiser of perception and multisensory integration more generally \citep{saltafossi_2025, kluger_2021, grund_2021}. Yet no published study has delivered PPS-task stimuli at controlled respiratory phases and fit the sigmoid under each. The single empirical study to embed a breathing manipulation in a PPS paradigm used breathing as an anxiety-induction procedure rather than as a phase manipulation \citep{spaccasassi_2019}, and the ``embreathment'' illusion demonstrates that breathing feedback on a seen avatar shifts corporeal awareness without engaging PPS measures directly \citep{monti2020embreathment, Allard2017, Betka2020, Betka2022}. The empirical gap is the simple intersection of these two literatures: respiratory phase has been shown to modulate generic multisensory integration, but not the body-referenced multisensory integration that defines PPS.

%% ---
%% PARAGRAPH 4b GOAL (one idea): PPS-plasticity manipulations to date
%% all operate via causal-inference-driven synchrony detection.
%% The missing manipulation is the exteroceptive-interoceptive
%% analogue: coupling a visible environmental signal to the breath.
%% Cardio-visual (Sel 2016) is the existing precedent; bioresponsive
%% VR makes the respiratory analogue feasible. This introduces
%% Viscereality as the paradigmatic answer without naming it yet.
%% ---
Addressing this gap invites a second, paradigmatic question. The plasticity manipulations that have most informed PPS theory (tool use, rubber hand and full-body illusions, audio--tactile far-and-near pairing) all work by introducing a tightly synchronous multimodal signal that the brain's causal-inference machinery treats as evidence of a shared generative source \citep{samad_2015, chancel_2021, costantini_2016, parise_2016}. An analogous manipulation has yet to be attempted for a purely visceral--exteroceptive pairing. The closest precedent is cardiovisual: \citet{sel_2016} showed that a pulsing visual shade locked to the observer's heartbeat increased identification with an otherwise ambiguous self--other face and modulated heartbeat-evoked potentials, with the effect scaling with individual interoceptive accuracy. Bioresponsive virtual reality makes the equivalent experiment for the respiratory channel technically feasible: visual elements of an immersive environment can be tied, in real time, to the participant's ongoing respiratory signal, an approach recently framed as \emph{exteroceptive--interoceptive sensory substitution} \citep{Goral2024}. In the taxonomy of PPS-plasticity paradigms this is a bioresponsive environmental-coupling manipulation, and it remains absent from the empirical record \citep{serino2017peripersonal, belli2021breathink, Belli2024}.

\subsection{The present study}

%% ---------------------------------------------------------------------
%% SUBSECTION 5 GOAL
%% What this subsection conveys: the design and its pre-registered
%% hypotheses. Each hypothesis should read as a direct consequence of
%% the argument above -- not as an independent claim. Order of
%% hypotheses follows the argumentative chain: coupling effect, then
%% phase-alone effect, then sharpness, then individual-differences
%% moderation.
%%
%% Two paragraphs, two ideas:
%%   (a) design summary -- phase-locked audio-tactile PPS under
%%       two within-subjects VR conditions (Viscereality vs Control).
%%   (b) four hypotheses, each mapped to the argument above.
%% ---------------------------------------------------------------------

%% ---
%% PARAGRAPH 5a GOAL (one idea): describe the design in enough detail
%% that the hypotheses in the next paragraph read as logical
%% consequences of the design. Emphasise (i) phase-locking of the
%% sigmoid to post-inhalation vs post-exhalation breath holds, and
%% (ii) the minimal-difference comparison between Viscereality (breath-
%% coupled orb) and Control (fixation ring on dark background).
%% ---
We address both gaps in a single within-subjects design. Participants perform an audio--tactile PPS task under controlled 4--4--4--4 box breathing (inhale, post-inhalation hold, exhale, post-exhalation hold). Audio--tactile stimuli are delivered only during the two breath-hold phases, so that each PPS sigmoid is fit separately for post-inhalation and post-exhalation holds, phase-locking the sigmoid to respiration for the first time in a PPS paradigm. Participants complete this task under two counterbalanced VR conditions. In the \emph{Viscereality} condition, an ovular orb in front of the participant expands and contracts in real time with their breathing, yielding a bioresponsive visual environment whose spatial dynamics are synchronous with the interoceptive breath signal. In the \emph{Control} condition, the identical task is performed with only a central fixation ring on a dark background, preserving attentional engagement and the underlying box-breathing rhythm but removing the breath-coupled spatial dynamics.

%% ---
%% PARAGRAPH 5b GOAL (one idea): state the four pre-registered
%% hypotheses. Each maps onto a prior subsection:
%%   H1 <- subsection 4b (coupling -> differential phase effect)
%%   H2 <- subsection 3b (phase alone -> pure excitability channel)
%%   H3 <- subsection 2 (softening / boundary dissolution)
%%   H4 <- subsection 2 (Ardizzi: IA moderates interoception-PPS link)
%% ---
Four hypotheses follow from the argument above. First, if coupling an exteroceptive spatial signal to the interoceptive breath signal engages the same causal-inference machinery that underlies body-ownership and PPS-plasticity effects \citep{samad_2015, chancel_2021, Ardizzi2018}, then the breath-coupled condition should produce a differential shift of the PPS boundary across respiratory phases, with PPS larger during post-inhalation holds (when breath-coupled visual space is expanded) than post-exhalation holds, yielding a Condition $\times$ Phase $\times$ SOA interaction. Second, comparing the Control condition alone tests whether respiratory phase, stripped of breath-coupled visuals, suffices to modulate the PPS sigmoid, arbitrating between phase-locked excitability accounts \citep{saltafossi_2025, kluger_2021, grund_2021} and coupling-specific accounts. Third, if Viscereality softens rather than simply translates the self--world boundary in the manner of contemplative and sensory-attenuation interventions \citep{Chiarella2024, hanley_2020, hruby_2024, nguyen_2020}, we expect a main effect of Condition on the sigmoid slope, with shallower boundaries under Viscereality than Control. Fourth, if the effects are mediated by how available the interoceptive signal is to perception, trait interoceptive awareness (MAIA-2) should moderate the magnitude of phase-locked modulation, with a stronger moderation under Viscereality where the internal signal is made externally salient \citep{Ardizzi2018, bogdanova_2021, Goral2024}. Pictographic measures of perceived body boundaries, spatial frame of reference, and small-self experience are collected as exploratory correlates of any Viscereality-induced boundary softening.
